\section{Conclusion}
\label{sec:conclusion}

Language models are deployed in conversations where users disagree with them, claim authority over them, press them across many turns, and describe stakes and threats. This survey defined that class of inputs as pressure, a property of the live prompt or conversational context that signals a performance demand and an implied consequence without stating the desired output, and distinguished it from the training-time factors that make a model pressure-prone, from jailbreaks that supply the disallowed output, and from persuasion that supplies an argument.

Under this definition the scattered literature resolves into five effects with a common cause. A model can be moved off a correct answer, into validating a stance it does not hold, into asserting something it does not believe or hiding something it did, into confidence and reasoning that no longer track correctness, and into weakened refusals and self-protective action. Each has been documented on its own, under its own name, by a community that rarely cites the others. Set side by side they correlate across models, chain within a single conversation, and respond to the same manipulation types, which is the case for treating them as one phenomenon.

Three findings only became visible once pressure was the unit of analysis. Capability does not buy robustness monotonically, and what it does buy is resistance on questions the model can verify rather than on the judgment questions where the stakes are highest. Models under the same pressure fail in different ways depending on who trained them, differences that are linearly encoded and switchable. And every mitigation tested so far reduces one failure at the cost of another, so the practical target is not eliminating the effect but choosing which trade-off to accept, a choice that a formal notion of pressure at least makes statable. The evidence is uneven: the cleanest results concern epistemic instability on verifiable items, and the most dramatic, on self-preservation, rest on contrived scenarios and unreplicated reports that we have flagged rather than relied on. What the field has not shown, despite repeated claims, is that pressure reliably helps. What it has shown, consistently, is that pressure is enough on its own to move a model off a correct answer, an honest statement, or a safe refusal, and that is reason enough to treat it as a variable worth naming.
